\section{Exact hardness at width two}\label{sec:hardness}
The bound of Theorem~\ref{thm:main} is polynomial in the values of the
numerical bounds $C,H,\mu^{-1},R,\sigma^{-1}$, not in their bit lengths,
so it is exponential in the input length when, for example, $C$ is
exponentially large. This section shows that some numerical dependence
of this kind is necessary. Theorem~\ref{thm:hardness} proves that exact
optimization is NP-complete already for matrices of bandwidth two that
are arbitrarily close to the identity. Bandwidth one is different: the
support graph is then a union of paths, and the path algorithm of
\cite{LiuPath} and the tree algorithm of \cite{BhathenaTrees} solve the
problem with quadratically many arithmetic operations.

Section~\ref{sec:hardness-perturbed} carries the hardness over to the
perturbed problem (Corollary~\ref{cor:magnitude}), and
Section~\ref{sec:related} compares both results with earlier lower
bounds.

\subsection{Idea of the reduction}\label{sec:hardness-idea}
All reductions start from SUBSET SUM \cite{Karp}: given positive
integers $a_1,\ldots,a_n$ and $B$, decide whether $\sum_ia_iz_i=B$ for
some $z\in\{0,1\}^n$. In the indicator problem we build, the indicators
choose the subset, and a chain of continuous variables adds up the
chosen numbers.

Each number $a_i$ gets a continuous \emph{control} $x_i$ with indicator
$z_i$, and the objective contains the term $x_i^2-2A_ix_i+A_i^2z_i$,
where the target $A_i$ is a scaled copy of $a_i$. If $z_i=0$, then
$x_i=0$ and the term is zero. If $z_i=1$, the term equals
$(x_i-A_i)^2$, which is zero at $x_i=A_i$. Both choices therefore cost
nothing as long as $x_i=A_iz_i$, so the indicators may select any
subset and the controls record it.

A chain of continuous \emph{states} $s_1,\ldots,s_n$, with $s_0=0$,
adds up the recorded numbers. For a small rational $\theta>0$, the
objective contains the squared \emph{dynamic residuals}
$(s_i-\theta s_{i-1}-\theta x_i)^2$. When all of them vanish, the states
follow the recursion $s_i=\theta s_{i-1}+\theta x_i$, and unrolling it
gives $s_n=\sum_i\theta^{n-i+1}x_i$. Each step shrinks the earlier part
of the sum by the factor $\theta$, and the targets undo this damping:
with $A_i=a_i\theta^{-(n-i+1)}$ and $x_i=A_iz_i$, the last state is
$s_n=\sum_ia_iz_i$. A final term $\theta^2(s_n-B)^2$ compares this sum
with $B$. On a YES instance, a solution of SUBSET SUM makes every square
vanish.

The matrix is close to the identity when $\theta$ is small. Each
variable appears with coefficient one in exactly one square, the control
$x_i$ in $x_i^2$ and the state $s_i$ in its own residual, and every other
square that contains it multiplies it by $\theta$. Each residual involves
only $s_{i-1}$, $x_i$, and $s_i$, which are at most two positions apart
in the order $(x_1,s_1,\ldots,x_n,s_n)$, so the matrix has bandwidth
two. Because the coupling is small, the targets $A_i$, and with them
the linear coefficients and the penalties, grow like powers of
$1/\theta$.

The gap between YES and NO instances comes from integrality. On a NO
instance, $B-\sum_ia_iz_i$ is a nonzero integer for every $z$, so its
absolute value is at least one. Unrolling the recursion writes this
integer as a fixed linear combination of the residuals: the control
residuals $x_i-A_iz_i$, the dynamic residuals, and $s_n-B$. The
residuals therefore cannot all be small, and the Cauchy--Schwarz
inequality turns this into a lower bound on their cost in the objective.
The bound is the reciprocal of a sum of squared weights, which are powers
of $\theta$, and this sum is bounded independently of $n$. Every NO
solution therefore pays more than a constant $\Delta>0$ that depends
only on $\theta$. Every continuous variable in \eqref{eq:problem} has
its own indicator, so the states need indicators $\zeta_i$ as well.
They get the small penalty $\Delta/(4n)$, so switching all of them on
costs only $\Delta/4$, well below the gap.

The construction just described is case~(1) of
Theorem~\ref{thm:hardness}. It has a fixed matrix and a constant gap,
but exponentially large coefficients and penalties. Case~(2) rescales
the variables and the objective of case~(1) by a common factor, which
bounds the data but can make the gap exponentially small. Case~(3)
moves the numbers $a_i$ from the targets into the coefficients of the
chain. The matrix then depends on the input, and every target and every
penalty equals one. A state indicator now costs as much as a control
indicator, and since the states are small, switching one off would cost
little in the residuals and save a whole unit of penalty. Shifting all
states by a constant $D$ prevents this, because an inactive state must
be zero, far from its required value near $D$.

\subsection{Statement and proof}
We prove hardness for the \emph{rational threshold decision problem}:
given a rational instance of \eqref{eq:problem} with $\xi=0$ and a
rational $\kappa$, decide whether the optimal value is at most $\kappa$.
The objectives below contain additive constants, which let us write the
residuals as squares. Subtracting such a constant from both the
objective and $\kappa$ gives an instance of \eqref{eq:problem}.

\begin{theorem}\label{thm:hardness}
Fix a rational $0<\theta\le1/10$, and put
\[
 \Delta=\frac{\theta^2(1-\theta^2)}{1+\theta^4}.
\]
The rational threshold decision problem for indicator quadratic optimization
is NP-complete under each of the following restrictions. In every case the
matrix has bandwidth two, maximum graph degree at most four, treewidth and
pathwidth two, and
\begin{equation}\label{eq:hard-spectrum}
 (1-3\theta)I\preceq Q\preceq(1+3\theta+2\theta^2)I.
\end{equation}
\begin{enumerate}
\item[(1)] The matrix depends only on the dimension and $\theta$, and
all penalties are positive. The reduction produces instances with
optimal value at most $\Delta/4$ from YES instances and greater than
$\Delta$ from NO instances. Linear coefficients and penalties may have
exponentially large magnitudes.
\item[(2)] The same matrix family, with $|c_i|\le2$ and
$0<\lambda_i\le1$. The gap between YES and NO values may be
exponentially small.
\item[(3)] All penalties equal one and $|c_i|\le9$. The matrix depends
on the input, through nonzero entries that may be exponentially small.
\end{enumerate}
In all cases the data have polynomial binary encoding length.
\end{theorem}

The three cases bound different numerical quantities:
\begin{center}
\small
\begin{tabular}{@{}lllll@{}}
Case & Matrix $Q$ & $|c_i|$ & Penalties $\lambda_i$
     & Gap between YES and NO values\\
\hline
(1) & fixed & exponential & positive, exponential
    & constant ($\le\Delta/4$ vs.\ $>\Delta$)\\
(2) & fixed & $\le2$ & in $(0,1]$ & may be exponentially small\\
(3) & input-dependent & $\le9$ & all equal to $1$
    & may be exponentially small
\end{tabular}
\end{center}
Here ``fixed'' means that $Q$ depends only on the dimension and
$\theta$, and ``exponential'' means possibly exponentially large in the
dimension. Corollary~\ref{cor:magnitude} builds on case~(1). In
case~(2) the optimal solutions are also bounded (see the proof), so only
the gap is small. Deciding such an instance through the certificate
\eqref{eq:certgap} needs $n\sigma$ below the gap, so $\sigma^{-1}$, and
with it the bound of Theorem~\ref{thm:main}, can be exponential.
Case~(3) shows that hardness does not need varied penalties, although no
case combines an input-independent matrix with unit penalties. Because
the reductions start from SUBSET SUM with binary-encoded numbers, they
do not show strong NP-hardness.

\begin{proof}
Take a SUBSET SUM instance with positive integers $a_1,\ldots,a_n,B$,
where $n\ge2$. As above, the variables $x_i$ are controls and the
variables $s_i$ (in case~(3), $y_i$) are states.

\emph{Case (1).} Define $A_i=a_i\theta^{-(n-i+1)}$, $\eta=\Delta/(4n)$,
and $s_0=0$. Use indicators $z_i$ for $x_i$ and $\zeta_i$ for $s_i$, and
minimize
\begin{equation}\label{eq:hard-objective}
 \Phi=\sum_i(x_i^2-2A_ix_i+A_i^2z_i)
   +\sum_i(s_i-\theta s_{i-1}-\theta x_i)^2
   +\theta^2(s_n-B)^2+\eta\sum_i\zeta_i.
\end{equation}
The penalties are $A_i^2$ for $z_i$ and $\eta$ for $\zeta_i$, and the
quadratic part of $\Phi$ does not depend on the indicators. On
indicator-feasible points $x_i=x_iz_i$ and $z_i=z_i^2$, so the first sum
equals $\sum_i(x_i-A_iz_i)^2$. If $\sum_i a_i z_i=B$, take
$x_i=A_i z_i$, $s_i=\theta s_{i-1}+\theta x_i$, and every $\zeta_i=1$.
Then $s_n=\sum_i\theta^{n-i+1}A_iz_i=B$, all squares vanish, and
$\Phi=n\eta=\Delta/4$.

For the converse, define the control residuals $e_i=x_i-A_i z_i$, the
dynamic residuals $r_i=s_i-\theta s_{i-1}-\theta x_i$, and $u=s_n-B$.
Unrolling the recursion for $s_n$ shows that every feasible point
satisfies
\[
 B-\sum_i a_i z_i
 =\sum_i\theta^{n-i+1}e_i+\sum_i\theta^{n-i}r_i-u.
\]
On a NO instance the left side has absolute value at least one, and
Cauchy--Schwarz, applied to $(e,r,\theta u)$, gives
\[
 \Phi\ge\|e\|^2+\|r\|^2+\theta^2u^2
 \ge\left[\theta^{-2}+(1+\theta^2)
                    \sum_{j=0}^{n-1}\theta^{2j}\right]^{-1}>\Delta.
\]
The last inequality uses $\sum_{j=0}^{n-1}\theta^{2j}<1/(1-\theta^2)$
and $1/\Delta=\theta^{-2}+(1+\theta^2)/(1-\theta^2)$.

In the order $(x_1,s_1,\ldots,x_n,s_n)$, every diagonal entry is
$1+\theta^2$, and the only off-diagonal entries are
\begin{equation}\label{eq:hard-edges}
 Q_{x_i,s_i}=-\theta,\qquad
 Q_{s_{i-1},s_i}=-\theta,\qquad
 Q_{s_{i-1},x_i}=\theta^2\quad(i\ge2),
\end{equation}
with only the first type present at $i=1$. The bags
$\{x_1,s_1\},\{s_1,x_2,s_2\},\ldots,\{s_{n-1},x_n,s_n\}$
give a width-two path decomposition, and since $n\ge2$, the triangle
$\{s_1,x_2,s_2\}$ makes the width exactly two. The maximum degree is at
most four. Absolute off-diagonal row sums are at most $3\theta+\theta^2$,
so Gershgorin's theorem proves \eqref{eq:hard-spectrum}; in particular,
the matrix is strictly diagonally dominant. For fixed rational $\theta$,
the powers defining $A_i$ and all remaining data have polynomial
encoding length, so the rational threshold $\Delta/4$ proves
NP-hardness. The gap also rules out approximation: an estimate within
absolute error $\Delta/4$ is at most $\Delta/2$ on YES instances and
greater than $3\Delta/4$ on NO instances. Hence approximating the
optimal value within absolute error $\Delta/4$ is also NP-hard.

\emph{Case (2).} Let $M=1+B+\sum_iA_i$, substitute
$(x,s)=M(\widehat x,\widehat s)$ in \eqref{eq:hard-objective}, and
divide the objective by $M^2$. The matrix is unchanged, $A_i$ and $B$
become $A_i/M$ and $B/M$, and the state penalty becomes $\eta/M^2$. This
gives the stated coefficient and penalty bounds and multiplies both
thresholds by $M^{-2}$. Optimal solutions are also bounded. The all-zero
point has value at most $\theta^2$, so at an optimum every control and
dynamic residual has absolute value at most $\theta$. Hence
$|\widehat x_i|\le1+\theta$ and
$|\widehat s_i|\le\theta(2+\theta)/(1-\theta)<1$.

\emph{Case (3).} We may assume $B\le A:=\sum_i a_i$, since otherwise the
answer is NO. Set
\[
 b_i=\frac{a_i\theta^i}{A},\qquad
 T=\frac{B\theta^n}{A},\qquad D=4,
\]
so $0<b_i\le\theta$ and $0<T\le1$. With $y_0=D$, indicators $z_i$ for
$x_i$ and $\zeta_i$ for $y_i$, and no constraints beyond the indicator
constraints, minimize
\begin{equation}\label{eq:unit-hard}
 G=\sum_i(x_i^2-2x_i+z_i)
 +\sum_i[y_i-\theta y_{i-1}-b_i x_i-D(1-\theta)]^2
 +\theta^2(y_n-D-T)^2+\sum_i\zeta_i.
\end{equation}
The first dynamic residual is $y_1-b_1x_1-D$. The shift $D$ makes
switching off a state indicator expensive, as shown below.
For each $z$, let $s_0^z=0$ and $s_i^z=\theta s_{i-1}^z+b_i z_i$; these
are the values of $y_i-D$ when $x=z$ and every dynamic residual
vanishes. Then $0\le s_i^z\le\theta^i$, and
$s_n^z=(\theta^n/A)\sum_ia_iz_i$ equals $T$ exactly when
$\sum_i a_i z_i=B$.
Put $e=x-z$ and $d=y-D\mathbf1-s^z$, and let $L$ be the lower shift
matrix. On feasible points, $x_i^2-2x_i+z_i=e_i^2$, and the $i$th
dynamic residual equals $d_i-\theta d_{i-1}-b_ie_i$ with $d_0=0$. The
control and dynamic squares therefore satisfy
\begin{equation}\label{eq:activation-quadratic}
 \|e\|^2+\|(I-\theta L)d-\operatorname{Diag}(b)e\|^2
 \ge(1-3\theta)(\|e\|^2+\|d\|^2).
\end{equation}
Indeed, in this quadratic form a control row has diagonal $1+b_i^2$
and off-diagonal absolute sum at most $(1+\theta)b_i$.
An internal state row has diagonal $1+\theta^2$ and row sum at most
$2\theta+b_i+\theta b_{i+1}\le3\theta+\theta^2$, and
the final row has diagonal one and row sum at most $2\theta$.
Gershgorin's theorem gives \eqref{eq:activation-quadratic}.

If $k$ state indicators vanish, the corresponding $d_i\le-D$. Since
$\theta\le1/10$ and $D=4$,
\[
 G\ge n-k+(1-3\theta)D^2k\ge n+10.2k.
\]
The feasible choice $x=z=0$, $y=D\mathbf1$, $\zeta=\mathbf1$ has
value $n+\theta^2T^2<n+1$. Consequently every optimum has all state
indicators on. Its value is $n$ plus a sum of squares, and it equals $n$
exactly for a YES instance. Positive definiteness gives attainment, so
every NO value is strictly greater than $n$.
To quantify the gap, let
$r_i=y_i-\theta y_{i-1}-b_ix_i-D(1-\theta)$ and $u=y_n-D-T$.
Using $e=x-z$ and $y_0=D$, the recurrence gives
\[
 \frac{\theta^n}{A}\left(B-\sum_i a_i z_i\right)
 =\sum_i\frac{\theta^n a_i}{A}e_i
       +\sum_i\theta^{n-i}r_i-u.
\]
At a NO optimum, Cauchy--Schwarz applied to $(e,r,\theta u)$ gives
\[
 G-n\ge\frac{\theta^{2n}/A^2}
 {\theta^{2n}\sum_i a_i^2/A^2+
       \sum_{j=0}^{n-1}\theta^{2j}+\theta^{-2}}>0.
\]
The numerator can be exponentially small.

The quadratic matrix of \eqref{eq:unit-hard} has diagonals $1+b_i^2$ on
controls and $1+\theta^2$ on states. Its off-diagonal entries are
$-b_i,-\theta,\theta b_i$ in the positions of \eqref{eq:hard-edges}.
All are nonzero, and the row bounds above prove the graph claims and
\eqref{eq:hard-spectrum}.
Writing $h_1=D$, $h_i=D(1-\theta)$ for $i>1$, the linear coefficients are
$-2+2h_i b_i$ on controls, $-2h_i+2\theta h_{i+1}$ on nonterminal states,
and $-2h_n-2\theta^2(D+T)$ on the last state. Their magnitudes are below
nine. The constant is at most $16n+1/4$, and all data have polynomial
encoding length.

\emph{Membership in NP.} For any rational positive-definite instance, a
support $J$ is a certificate. Its unique continuous optimum is
$-Q_{JJ}^{-1}c_J/2$, and its value has polynomial rational encoding
length by determinant bounds. Exact linear algebra verifies the
threshold in polynomial time. Supports containing zero optimizing
coordinates remain valid certificates. This proves membership in NP in
all three cases.
\end{proof}

All three cases come arbitrarily close to the identity. Dividing the
objective of cases~(1) and~(2) by $1+\theta^2$ gives a unit-diagonal
matrix with $\|Q-I\|_2\le(3\theta+\theta^2)/(1+\theta^2)\le4\theta$,
and case~(3) has $\|Q-I\|_2\le3\theta+2\theta^2$ by
\eqref{eq:hard-spectrum}. At $Q=I$ the problem separates into scalar
problems, so hardness appears arbitrarily close to a trivial case.

The instances also avoid a known polynomial case. Problem
\eqref{eq:problem} is polynomial-time solvable when $Q$ is a Stieltjes
matrix, that is, has nonpositive off-diagonal entries, and the linear
coefficients have one sign \cite{AtamturkGomezM}, as recalled by
Liu et al.~\cite{LiuStieltjes}. The linear coefficients in all three
cases are nonpositive, so it matters that the matrices are not
Stieltjes. Changing the signs of some continuous variables turns an
instance into an equivalent one, but it cannot make these matrices
Stieltjes. Each triangle of the support graph has a positive product of
off-diagonal entries, sign changes preserve this product, and in a
Stieltjes matrix the product would be negative.

Although we call the variables controls and states, the instances lie
outside the multi-period model of Lee, G\'omez, and Atamt\"urk
\cite{LeeMultiperiod}. There the dynamics are exact constraints and the
quadratic cost involves only the states; here the dynamics are only
penalized, and every control has its own square. Even with all state
indicators on, minimizing over the states leaves the control matrix
$I+\alpha ww^T$, where $\alpha>0$, $w_i=\theta^{n-i+1}$ in cases~(1)
and~(2), and $w_i=\theta^{n-i}b_i$ in case~(3). Its inverse has no zero
entries, so for $n\ge3$ it is not tridiagonal.

\subsection{Consequence for perturbed penalties}
\label{sec:hardness-perturbed}
Theorem~\ref{thm:hardness} concerns the unperturbed problem.
Corollary~\ref{cor:magnitude} below fixes the matrix family of
case~(1), hence $H$ and $\mu$, and the perturbation scale $\sigma$; it
does not involve $R$; and it scales only the linear data $c$ and
$\lambda$.

\begin{corollary}\label{cor:magnitude}
Fix a rational $\sigma>0$, and let each $\xi_i$ be drawn independently
and uniformly from the perturbation grid \eqref{eq:noisegrid}, with $N$
as in Theorem~\ref{thm:main}. Suppose an algorithm returns the exact
optimal value of \eqref{eq:problem} for every draw, and there is a
polynomial $p$ such that, on every instance of encoding length $\ell$,
its expected bit complexity over the draw is at most $p(\ell)$. Then
$\mathrm{NP}\subseteq\mathrm{ZPP}$. This holds even if the algorithm is
required to work only for the fixed matrix family of
Theorem~\ref{thm:hardness}(1).
\end{corollary}
The proof is a scaling argument. Multiplying the linear coefficients of
a case-(1) instance by $\gamma$ and its penalties by $\gamma^2$
multiplies the gap between YES and NO values by $\gamma^2$,
while the perturbation changes every value by at most $2n\sigma$,
independently of $\gamma$. For large $\gamma$ the exact perturbed value
therefore still decides the SUBSET SUM instance. Beier and V\"ocking use
the same scaling to show that polynomial smoothed complexity implies a
randomized pseudopolynomial algorithm
\cite[conference version, Theorem~3]{BeierVocking}; see also
\cite[Theorem~6.2]{RoglinTeng}.
\begin{proof}
Start from a case-(1) instance. It has $n_{\mathrm{tot}}=2n$ indicators,
and its smallest penalty is $\eta=\Delta/(4n)$. Choose an integer
$\gamma\ge1$ with $3\gamma^2\Delta/4>2n_{\mathrm{tot}}\sigma$ and
$\gamma^2\Delta/(4n)>\sigma$. Such a $\gamma$ has polynomial encoding
length and is computable by rational comparisons.
Let $c_0$ denote the additive constant in the expanded objective
\eqref{eq:hard-objective}. Replace $(c,\lambda,c_0)$ by
$(\gamma c,\gamma^2\lambda,\gamma^2c_0)$, keeping $Q$.
Substituting $\gamma x'$ for the vector $x$ of all continuous variables
shows that every feasible value is multiplied by $\gamma^2$. A
perturbation with $|\xi_i|\le\sigma$ changes the optimum by at most
$n_{\mathrm{tot}}\sigma$. The perturbed YES value is therefore at most
$\gamma^2\Delta/4+n_{\mathrm{tot}}\sigma$, whereas the NO value exceeds
$\gamma^2\Delta-n_{\mathrm{tot}}\sigma$.
These ranges are disjoint for every draw, and all perturbed penalties
remain positive. Sample the perturbation, run the solver, and compare
its exact answer with the midpoint of these two bounds. The assumed
expected bound gives a zero-error expected-polynomial algorithm for
SUBSET SUM.
\end{proof}
The scaled case-(1) instances have exponentially large $C$, so the
bound of Theorem~\ref{thm:main} is exponential for them. The argument
scales the penalties together with the linear coefficients, so it gives
no obstruction when the penalties are bounded.
